\documentclass[11pt,a4paper]{article}
\usepackage{fontspec,amsmath,amssymb,amsthm,microtype}
\setmainfont{lmroman10-regular.otf}[BoldFont=lmroman10-bold.otf,ItalicFont=lmroman10-italic.otf,BoldItalicFont=lmroman10-bolditalic.otf]
\usepackage{xcolor}
\usepackage[most]{tcolorbox}
\usepackage{mathrsfs}
\providecommand{\lt}{<}
\providecommand{\gt}{>}
\setmonofont{DejaVuSansMono.ttf}[BoldFont=DejaVuSansMono-Bold.ttf,ItalicFont=DejaVuSansMono-Oblique.ttf,BoldItalicFont=DejaVuSansMono-BoldOblique.ttf,Scale=0.86]
\newfontfamily\leanfallback{STIXTwoMath-Regular.otf}[Scale=0.86]
\definecolor{leanrule}{gray}{0.62}
\definecolor{leanback}{gray}{0.962}
\definecolor{leantitle}{gray}{0.22}
\definecolor{leanmuted}{gray}{0.38}
\newtcolorbox{leanbox}[1]{enhanced,breakable,colback=leanback,colframe=leanrule,
  boxrule=0.4pt,arc=0.6mm,left=2.4mm,right=2.4mm,top=1.8mm,bottom=1.8mm,
  before skip=0.8em,after skip=1em,fonttitle=\small\bfseries,coltitle=leantitle,
  colbacktitle=leanback,titlerule=0.3pt,title={#1}}
\newlength{\leancharwidth}
\newenvironment{leancode}{\par\smallskip\begingroup\ttfamily\footnotesize\raggedright
  \settowidth{\leancharwidth}{M}\setlength{\parskip}{0pt}\setlength{\parindent}{0pt}
  \frenchspacing}{\par\endgroup\smallskip}
\newcommand{\leanline}[2]{\hangindent=\dimexpr#1\leancharwidth+4\leancharwidth\relax
  \hangafter=1\noindent\hspace*{#1\leancharwidth}#2\par}
\newcommand{\leansource}[1]{\par{\footnotesize\color{leanmuted}#1}\par}
\newcommand{\leanlabel}[1]{\par\smallskip\noindent{\small\bfseries #1}\par}
\usepackage[margin=25mm]{geometry}
\usepackage[unicode,hidelinks]{hyperref}
\hypersetup{pdfauthor={LangevinLean Project}}
\setlength{\emergencystretch}{3em}
\setlength{\parskip}{0.4em}
\setlength{\parindent}{0pt}
\allowdisplaybreaks[2]
\raggedbottom
\title{\texorpdfstring{Couplings and the quadratic transport distance}{Couplings and the quadratic transport distance}}
\author{LangevinLean Project}
\date{}
\begin{document}\maketitle
Draft mathematical exposition; these notes may contain mistakes.


Let \(H\) be a normed additive group (for instance a normed space) whose topology is second countable, equipped with its Borel \(\sigma\)-algebra; then the Borel \(\sigma\)-algebra of \(H\times H\) is the product \(\sigma\)-algebra and \((x,y)\mapsto x-y\) is measurable. For a measure \(\mu\) and a measurable map \(f\), the image measure is \(f_\#\mu(A)=\mu(f^{-1}(A))\). A \emph{coupling} of two measures \(\mu,\nu\) on \(H\) is a measure \(\gamma\) on \(H\times H\) whose marginals \((\mathrm{pr}_1)_\#\gamma\) and \((\mathrm{pr}_2)_\#\gamma\) are \(\mu\) and \(\nu\); the set of couplings is \(\Gamma(\mu,\nu)\).




For a measure \(P\) on \(\Omega\) and a function \(f:\Omega\to H\) we use the extended \(L^2\)-norm \(\|f\|_{2,P}=(\int\|f\|^2\,dP)^{1/2}\in[0,\infty]\) when \(f\) is \(P\)-almost everywhere strongly measurable (equal almost everywhere to a pointwise limit of measurable functions with finitely many values), and \(\|f\|_{2,P}=\infty\) otherwise. With this convention, Minkowski's inequality \(\|f-g\|_{2,P}\le\|f\|_{2,P}+\|g\|_{2,P}\) holds for all \(f,g\), and \(\|f\|_{2,P}\) does not change if \(f\) is modified on a null set. In \([0,\infty]\) we use \(a+\infty=\infty\), \(a\cdot\infty=\infty\) for \(a>0\) and \(0\cdot\infty=0\).




The quadratic cost of \(\gamma\) is \(c(\gamma)=\|(x,y)\mapsto x-y\|_{2,\gamma}=(\int\|x-y\|^2\,\gamma(dx,dy))^{1/2}\), and the quadratic Wasserstein distance is


\[W_2(\mu,\nu)=\inf_{\gamma\in\Gamma(\mu,\nu)}c(\gamma)\in[0,\infty],\]


with \(\inf\varnothing=\infty\). We also use the real number \(W_2^{\mathbb R}(\mu,\nu)\), equal to \(W_2(\mu,\nu)\) if this is finite and to \(0\) otherwise; every bound below on \(W_2^{\mathbb R}\) is accompanied by a proof or a hypothesis of finiteness. No optimal coupling and no metric property of \(W_2\) is used.




The main facts are: every coupling bounds \(W_2\) from above; the joint law \((X,Y)_\#P\) of two random vectors is a coupling of their laws, with cost \(\|X-Y\|_{2,P}\), so \(W_2(X_\#P,Y_\#P)\le\|X-Y\|_{2,P}\); and if \(X'-Y'=A-V\) with \(\|A\|\le q\|X-Y\|\) pointwise, then Minkowski's inequality gives \(\|X'-Y'\|_{2,P}\le q\|X-Y\|_{2,P}+\|V\|_{2,P}\). Combined, these turn a pointwise contraction plus a defect into a bound on the transport distance between the laws of the next pair. No independence between the compared random vectors is required.


\section{\texorpdfstring{Couplings}{Couplings}}


\textbf{Definition.} Let \(\mu,\nu\) be measures on \(H\). A coupling of \(\mu\) and \(\nu\) is a measure \(\gamma\) on \(H\times H\) such that


\[(\mathrm{pr}_1)_\#\gamma=\mu\quad\text{and}\quad(\mathrm{pr}_2)_\#\gamma=\nu,\qquad\text{i.e.}\quad\gamma(A\times H)=\mu(A),\ \gamma(H\times B)=\nu(B)\]


for all Borel sets \(A,B\subseteq H\).




The dependence between the two coordinates is arbitrary. For probability measures the product \(\mu\otimes\nu\) is a coupling, so \(\Gamma(\mu,\nu)\) is nonempty.

\begin{leanbox}{Lean: \texttt{Coupling}}
\leansource{\texttt{Coupling} \quad UnadjustedLangevin/\allowbreak{}Coupling.lean, lines 31–34}
\begin{leancode}
\leanline{0}{structure\ Coupling\ (μ\ ν\ :\ Measure\ H)\ where}
\leanline{2}{joint\ :\ Measure\ (H\ ×\ H)}
\leanline{2}{fst\ :\ joint.map\ Prod.fst\ =\ μ}
\leanline{2}{snd\ :\ joint.map\ Prod.snd\ =\ ν}
\end{leancode}
\leanlabel{How the Lean proof follows the argument}
\texttt{structure Coupling (\ensuremath{\mu} \ensuremath{\nu} : Measure H)} has three fields: \texttt{joint : Measure (H \ensuremath{\times} H)} is \(\gamma\); \texttt{fst : joint.\allowbreak{}map Prod.\allowbreak{}fst = \ensuremath{\mu}} and \texttt{snd : joint.\allowbreak{}map Prod.\allowbreak{}snd = \ensuremath{\nu}} are the two marginal conditions, where \texttt{\ensuremath{\gamma}.\allowbreak{}map f} is the image measure \(f_\#\gamma\) and \texttt{Prod.\allowbreak{}fst}, \texttt{Prod.\allowbreak{}snd} are the projections. An element of the structure is a coupling together with the proofs of its marginal conditions.
\end{leanbox}

\section{\texorpdfstring{The quadratic cost of a coupling}{The quadratic cost of a coupling}}


\textbf{Definition.} For a measure \(\gamma\) on \(H\times H\),


\[c(\gamma)=\|(x,y)\mapsto x-y\|_{2,\gamma}=\Bigl(\int_{H\times H}\|x-y\|^2\,\gamma(dx,dy)\Bigr)^{1/2}\in[0,\infty].\]


The displacement \((x,y)\mapsto x-y\) is continuous, hence measurable, so the integral formula applies. The value \(\infty\) is allowed; it occurs exactly when the second moment of the displacement is infinite.

\begin{leanbox}{Lean: \texttt{couplingCost}}
\leansource{\texttt{couplingCost} \quad UnadjustedLangevin/\allowbreak{}Coupling.lean, lines 37–38}
\begin{leancode}
\leanline{0}{def\ couplingCost\ (γ\ :\ Measure\ (H\ ×\ H))\ :\ ℝ≥0∞\ :=}
\leanline{2}{eLpNorm\ (fun\ p\ :\ H\ ×\ H\ ↦\ p.1\ -\ p.2)\ 2\ γ}
\end{leancode}
\leanlabel{How the Lean proof follows the argument}
\texttt{coupling\allowbreak{}Cost \ensuremath{\gamma}} is \texttt{e\allowbreak{}Lp\allowbreak{}Norm (fun p : H \ensuremath{\times} H \ensuremath{\mapsto} p.\allowbreak{}1 -\allowbreak{} p.\allowbreak{}2) 2 \ensuremath{\gamma}}. Mathlib's \texttt{e\allowbreak{}Lp\allowbreak{}Norm f p \ensuremath{\mu}} is the extended \(L^p\)-seminorm with values in \texttt{\ensuremath{\mathbb R}\ensuremath{\ge}0\ensuremath{\infty}}, the extended nonnegative reals \([0,\infty]\): for \(p=2\) and an almost everywhere strongly measurable \(f\) it equals \((\int\|f\|^2\,d\mu)^{1/2}\), and it is \(\infty\) for other functions. Keeping the value in \texttt{\ensuremath{\mathbb R}\ensuremath{\ge}0\ensuremath{\infty}} prevents an infinite cost from being silently replaced by a real number.
\end{leanbox}

\section{\texorpdfstring{The quadratic Wasserstein distance}{The quadratic Wasserstein distance}}


\textbf{Definition.} For measures \(\mu,\nu\) on \(H\),


\[W_2(\mu,\nu)=\inf_{\gamma\in\Gamma(\mu,\nu)}c(\gamma)\in[0,\infty],\]


the infimum being taken in \([0,\infty]\), so that \(W_2(\mu,\nu)=\infty\) if \(\Gamma(\mu,\nu)=\varnothing\) or every coupling has infinite cost.

\begin{leanbox}{Lean: \texttt{wasserstein2}}
\leansource{\texttt{wasserstein2} \quad UnadjustedLangevin/\allowbreak{}Coupling.lean, lines 41–42}
\begin{leancode}
\leanline{0}{def\ wasserstein2\ (μ\ ν\ :\ Measure\ H)\ :\ ℝ≥0∞\ :=}
\leanline{2}{{\leanfallback ⨅}\ c\ :\ Coupling\ μ\ ν,\ couplingCost\ c.joint}
\end{leancode}
\leanlabel{How the Lean proof follows the argument}
\texttt{wasserstein2 \ensuremath{\mu} \ensuremath{\nu}} is \texttt{{\leanfallback ⨅} c : Coupling \ensuremath{\mu} \ensuremath{\nu}, coupling\allowbreak{}Cost c.\allowbreak{}joint}, an indexed infimum in the complete lattice \texttt{\ensuremath{\mathbb R}\ensuremath{\ge}0\ensuremath{\infty}}; the infimum over an empty index type is the top element \(\infty\).
\end{leanbox}

\section{\texorpdfstring{The real-valued distance}{The real-valued distance}}


\textbf{Definition.} For measures \(\mu,\nu\) on \(H\), let \(W_2^{\mathbb R}(\mu,\nu)=W_2(\mu,\nu)\) if \(W_2(\mu,\nu)<\infty\), and \(W_2^{\mathbb R}(\mu,\nu)=0\) if \(W_2(\mu,\nu)=\infty\).




This real number coincides with the transport distance whenever the latter is finite. The theorems below that bound \(W_2^{\mathbb R}\) either assume finiteness or prove it from a coupling of finite cost, so the value \(0\) assigned to an infinite distance never enters a conclusion.

\begin{leanbox}{Lean: \texttt{wasserstein2Real}}
\leansource{\texttt{wasserstein2Real} \quad UnadjustedLangevin/\allowbreak{}Coupling.lean, lines 46–46}
\begin{leancode}
\leanline{0}{def\ wasserstein2Real\ (μ\ ν\ :\ Measure\ H)\ :\ ℝ\ :=\ (wasserstein2\ μ\ ν).toReal}
\end{leancode}
\leanlabel{How the Lean proof follows the argument}
\texttt{wasserstein2\allowbreak{}Real \ensuremath{\mu} \ensuremath{\nu}} is \texttt{(wasserstein2 \ensuremath{\mu} \ensuremath{\nu}).\allowbreak{}to\allowbreak{}Real}. The conversion \texttt{ENNReal.\allowbreak{}to\allowbreak{}Real : \ensuremath{\mathbb R}\ensuremath{\ge}0\ensuremath{\infty} \ensuremath{\to} \ensuremath{\mathbb R}} is the identity on finite values and maps \(\infty\) to \(0\), which is the convention of the definition.
\end{leanbox}

\section{\texorpdfstring{Every coupling gives an upper bound}{Every coupling gives an upper bound}}


\textbf{Claim.} For measures \(\mu,\nu\) on \(H\) and every \(\gamma\in\Gamma(\mu,\nu)\), \(W_2(\mu,\nu)\le c(\gamma)\).




\textbf{Proof.} An infimum is a lower bound of the family over which it is taken. \(\square\)

\begin{leanbox}{Lean: \texttt{wasserstein2\_\allowbreak{}le\_\allowbreak{}couplingCost}}
\leansource{\texttt{wasserstein2\_\allowbreak{}le\_\allowbreak{}couplingCost} \quad UnadjustedLangevin/\allowbreak{}Coupling.lean, lines 50–52}
\begin{leancode}
\leanline{0}{theorem\ wasserstein2\_le\_couplingCost\ \{μ\ ν\ :\ Measure\ H\}\ (c\ :\ Coupling\ μ\ ν)\ :}
\leanline{4}{wasserstein2\ μ\ ν\ ≤\ couplingCost\ c.joint}
\end{leancode}
\leanlabel{How the Lean proof follows the argument}
\texttt{wasserstein2\_\allowbreak{}le\_\allowbreak{}coupling\allowbreak{}Cost} is \texttt{i\allowbreak{}Inf\_\allowbreak{}le \_\allowbreak{} c}. The line \texttt{omit [Borel\allowbreak{}Space H] [Second\allowbreak{}Countable\allowbreak{}Topology H] in} records that these two assumptions are not needed.
\end{leanbox}

\section{\texorpdfstring{The joint law of two random vectors}{The joint law of two random vectors}}


\textbf{Claim.} Let \(P\) be a measure on \(\Omega\) and \(X,Y:\Omega\to H\) measurable. Then the joint law \(\gamma=(X,Y)_\#P\), the image of \(P\) under \(\omega\mapsto(X(\omega),Y(\omega))\), is a coupling of \(X_\#P\) and \(Y_\#P\).




\textbf{Proof.} The map \((X,Y):\Omega\to H\times H\) is measurable for the product \(\sigma\)-algebra because both coordinates are measurable, and the projections are measurable. For measurable maps the image of an image measure is the image under the composition: \((\mathrm{pr}_1)_\#(X,Y)_\#P=(\mathrm{pr}_1\circ(X,Y))_\#P=X_\#P\), since \(\mathrm{pr}_1\circ(X,Y)=X\). Likewise \((\mathrm{pr}_2)_\#\gamma=Y_\#P\). \(\square\)




No relation between \(X\) and \(Y\) (such as independence) is required.

\begin{leanbox}{Lean: \texttt{randomVariableCoupling}}
\leansource{\texttt{randomVariableCoupling} \quad UnadjustedLangevin/\allowbreak{}Coupling.lean, lines 57–65}
\begin{leancode}
\leanline{0}{def\ randomVariableCoupling\ (X\ Y\ :\ Ω\ →\ H)\ (hX\ :\ Measurable\ X)\ (hY\ :\ Measurable\ Y)\ :}
\leanline{4}{Coupling\ (P.map\ X)\ (P.map\ Y)\ where}
\leanline{2}{joint\ :=\ P.map\ (fun\ ω\ ↦\ (X\ ω,\ Y\ ω))}
\leanline{2}{fst\ :=\ by}
\leanline{4}{rw\ [Measure.map\_map\ measurable\_fst\ (hX.prodMk\ hY)]}
\leanline{4}{rfl}
\leanline{2}{snd\ :=\ by}
\leanline{4}{rw\ [Measure.map\_map\ measurable\_snd\ (hX.prodMk\ hY)]}
\leanline{4}{rfl}
\end{leancode}
\leanlabel{How the Lean proof follows the argument}
\texttt{random\allowbreak{}Variable\allowbreak{}Coupling P X Y h\allowbreak{}X h\allowbreak{}Y} constructs the coupling: \texttt{joint := P.\allowbreak{}map (fun \ensuremath{\omega} \ensuremath{\mapsto} (X \ensuremath{\omega}, Y \ensuremath{\omega}))}, and each marginal condition is proved by \texttt{rw [Measure.\allowbreak{}map\_\allowbreak{}map measurable\_\allowbreak{}fst (h\allowbreak{}X.\allowbreak{}prod\allowbreak{}Mk h\allowbreak{}Y)]} (Mathlib's composition rule for image measures, with measurability of the pair from \texttt{Measurable.\allowbreak{}prod\allowbreak{}Mk}) followed by \texttt{rfl}.
\end{leanbox}

\section{\texorpdfstring{Joint laws of almost everywhere measurable vectors}{Joint laws of almost everywhere measurable vectors}}


\textbf{Claim.} Let \(P\) be a measure on \(\Omega\) and \(X,Y:\Omega\to H\) be \(P\)-almost everywhere measurable. Then \((X,Y)_\#P\) is a coupling of \(X_\#P\) and \(Y_\#P\).




\textbf{Proof.} For an almost everywhere measurable map, the image measure is computed with any measurable map equal to it almost everywhere, and does not depend on that choice. The pair \((X,Y)\) is almost everywhere measurable, the projections are measurable, and the composition rule \((\mathrm{pr}_i)_\#(X,Y)_\#P=(\mathrm{pr}_i\circ(X,Y))_\#P\) holds in this setting; the proof is then that of the previous step. \(\square\)




This form applies to representatives of classes in \(L^2\), which are in general only almost everywhere measurable.

\begin{leanbox}{Lean: \texttt{aeRandomVariableCoupling}}
\leansource{\texttt{aeRandomVariableCoupling} \quad UnadjustedLangevin/\allowbreak{}Coupling.lean, lines 68–77}
\begin{leancode}
\leanline{0}{def\ aeRandomVariableCoupling\ (X\ Y\ :\ Ω\ →\ H)}
\leanline{4}{(hX\ :\ AEMeasurable\ X\ P)\ (hY\ :\ AEMeasurable\ Y\ P)\ :}
\leanline{4}{Coupling\ (P.map\ X)\ (P.map\ Y)\ where}
\leanline{2}{joint\ :=\ P.map\ (fun\ ω\ ↦\ (X\ ω,\ Y\ ω))}
\leanline{2}{fst\ :=\ by}
\leanline{4}{rw\ [AEMeasurable.map\_map\_of\_aemeasurable\ measurable\_fst.aemeasurable\ (hX.prodMk\ hY)]}
\leanline{4}{rfl}
\leanline{2}{snd\ :=\ by}
\leanline{4}{rw\ [AEMeasurable.map\_map\_of\_aemeasurable\ measurable\_snd.aemeasurable\ (hX.prodMk\ hY)]}
\leanline{4}{rfl}
\end{leancode}
\leanlabel{How the Lean proof follows the argument}
\texttt{ae\allowbreak{}Random\allowbreak{}Variable\allowbreak{}Coupling} has the same \texttt{joint}; the marginal conditions use \texttt{AEMeasurable.\allowbreak{}map\_\allowbreak{}map\_\allowbreak{}of\_\allowbreak{}aemeasurable}, the composition rule for image measures under almost everywhere measurable maps, with \texttt{h\allowbreak{}X.\allowbreak{}prod\allowbreak{}Mk h\allowbreak{}Y} for the pair.
\end{leanbox}

\section{\texorpdfstring{The cost of a joint law}{The cost of a joint law}}


\textbf{Claim.} For measurable \(X,Y:\Omega\to H\),


\[c\bigl((X,Y)_\#P\bigr)=\|X-Y\|_{2,P}.\]


\textbf{Proof.} Let \(\varphi(x,y)=x-y\), a continuous map on \(H\times H\). By the change-of-variables formula for image measures, \(\int\|\varphi\|^2\,d\bigl((X,Y)_\#P\bigr)=\int\|\varphi\circ(X,Y)\|^2\,dP\), and \(\varphi\circ(X,Y)=X-Y\). Taking square roots gives the claim; both sides are infinite simultaneously. \(\square\)

\begin{leanbox}{Lean: \texttt{randomVariableCoupling\_\allowbreak{}cost}}
\leansource{\texttt{randomVariableCoupling\_\allowbreak{}cost} \quad UnadjustedLangevin/\allowbreak{}Coupling.lean, lines 80–86}
\begin{leancode}
\leanline{0}{theorem\ randomVariableCoupling\_cost\ (X\ Y\ :\ Ω\ →\ H)}
\leanline{4}{(hX\ :\ Measurable\ X)\ (hY\ :\ Measurable\ Y)\ :}
\leanline{4}{couplingCost\ (randomVariableCoupling\ P\ X\ Y\ hX\ hY).joint\ =}
\leanline{6}{eLpNorm\ (fun\ ω\ ↦\ X\ ω\ -\ Y\ ω)\ 2\ P}
\end{leancode}
\leanlabel{How the Lean proof follows the argument}
\texttt{random\allowbreak{}Variable\allowbreak{}Coupling\_\allowbreak{}cost} unfolds the two definitions and applies Mathlib's \texttt{e\allowbreak{}Lp\allowbreak{}Norm\_\allowbreak{}map\_\allowbreak{}measure}, \(\|g\|_{p,f_\#\mu}=\|g\circ f\|_{p,\mu}\), whose hypotheses (strong measurability of the displacement, proved by \texttt{fun\_\allowbreak{}prop}, and almost everywhere measurability of the pair) are supplied.
\end{leanbox}

\section{\texorpdfstring{A coupled pair bounds the transport distance}{A coupled pair bounds the transport distance}}


\textbf{Claim.} For measurable \(X,Y:\Omega\to H\),


\[W_2(X_\#P,Y_\#P)\le\|X-Y\|_{2,P}.\]


\textbf{Proof.} By Step 6, \((X,Y)_\#P\in\Gamma(X_\#P,Y_\#P)\). By Step 5, \(W_2(X_\#P,Y_\#P)\le c((X,Y)_\#P)\), and by Step 8 this cost equals \(\|X-Y\|_{2,P}\). \(\square\)




Thus any two random vectors defined on a common probability space, for instance two copies of an algorithm driven by the same noise, give an upper bound for the transport distance between their laws.

\begin{leanbox}{Lean: \texttt{wasserstein2\_\allowbreak{}le\_\allowbreak{}random\_\allowbreak{}cost}}
\leansource{\texttt{wasserstein2\_\allowbreak{}le\_\allowbreak{}random\_\allowbreak{}cost} \quad UnadjustedLangevin/\allowbreak{}Coupling.lean, lines 89–93}
\begin{leancode}
\leanline{0}{theorem\ wasserstein2\_le\_random\_cost\ (X\ Y\ :\ Ω\ →\ H)}
\leanline{4}{(hX\ :\ Measurable\ X)\ (hY\ :\ Measurable\ Y)\ :}
\leanline{4}{wasserstein2\ (P.map\ X)\ (P.map\ Y)\ ≤\ eLpNorm\ (fun\ ω\ ↦\ X\ ω\ -\ Y\ ω)\ 2\ P}
\end{leancode}
\leanlabel{How the Lean proof follows the argument}
\texttt{wasserstein2\_\allowbreak{}le\_\allowbreak{}random\_\allowbreak{}cost} is \texttt{wasserstein2\_\allowbreak{}le\_\allowbreak{}coupling\allowbreak{}Cost (random\allowbreak{}Variable\allowbreak{}Coupling P X Y h\allowbreak{}X h\allowbreak{}Y)} rewritten by \texttt{random\allowbreak{}Variable\allowbreak{}Coupling\_\allowbreak{}cost} (inside \texttt{simpa}).
\end{leanbox}

\section{\texorpdfstring{The same bound for almost everywhere measurable vectors}{The same bound for almost everywhere measurable vectors}}


\textbf{Claim.} For \(P\)-almost everywhere measurable \(X,Y:\Omega\to H\),


\[W_2(X_\#P,Y_\#P)\le\|X-Y\|_{2,P}.\]


\textbf{Proof.} By Step 7, \((X,Y)_\#P\) is a coupling, so \(W_2(X_\#P,Y_\#P)\le c((X,Y)_\#P)\) by Step 5. The change of variables of Step 8 holds for the almost everywhere measurable map \((X,Y)\) and gives \(c((X,Y)_\#P)=\|X-Y\|_{2,P}\). \(\square\)

\begin{leanbox}{Lean: \texttt{wasserstein2\_\allowbreak{}le\_\allowbreak{}ae\_\allowbreak{}random\_\allowbreak{}cost}}
\leansource{\texttt{wasserstein2\_\allowbreak{}le\_\allowbreak{}ae\_\allowbreak{}random\_\allowbreak{}cost} \quad UnadjustedLangevin/\allowbreak{}Coupling.lean, lines 96–100}
\begin{leancode}
\leanline{0}{theorem\ wasserstein2\_le\_ae\_random\_cost\ (X\ Y\ :\ Ω\ →\ H)}
\leanline{4}{(hX\ :\ AEMeasurable\ X\ P)\ (hY\ :\ AEMeasurable\ Y\ P)\ :}
\leanline{4}{wasserstein2\ (P.map\ X)\ (P.map\ Y)\ ≤\ eLpNorm\ (fun\ ω\ ↦\ X\ ω\ -\ Y\ ω)\ 2\ P}
\end{leancode}
\leanlabel{How the Lean proof follows the argument}
\texttt{wasserstein2\_\allowbreak{}le\_\allowbreak{}ae\_\allowbreak{}random\_\allowbreak{}cost} composes \texttt{wasserstein2\_\allowbreak{}le\_\allowbreak{}coupling\allowbreak{}Cost (ae\allowbreak{}Random\allowbreak{}Variable\allowbreak{}Coupling P X Y h\allowbreak{}X h\allowbreak{}Y)} with the identity \texttt{e\allowbreak{}Lp\allowbreak{}Norm\_\allowbreak{}map\_\allowbreak{}measure} (via \texttt{trans\_\allowbreak{}eq}).
\end{leanbox}

\section{\texorpdfstring{A finite coupled cost gives a real bound}{A finite coupled cost gives a real bound}}


\textbf{Claim.} Let \(X,Y:\Omega\to H\) be measurable with \(X-Y\) square integrable, i.e. \(\|X-Y\|_{2,P}<\infty\). Then \(W_2(X_\#P,Y_\#P)<\infty\) and


\[W_2^{\mathbb R}(X_\#P,Y_\#P)\le\|X-Y\|_{2,P}.\]


\textbf{Proof.} By Step 9, \(W_2(X_\#P,Y_\#P)\le\|X-Y\|_{2,P}<\infty\). Both sides are finite, so the same inequality holds between real numbers, and the left-hand side is \(W_2^{\mathbb R}(X_\#P,Y_\#P)\). \(\square\)




Only the difference \(X-Y\) needs a finite second moment, not \(X\) and \(Y\) separately.

\begin{leanbox}{Lean: \texttt{wasserstein2\_\allowbreak{}le\_\allowbreak{}random\_\allowbreak{}cost\_\allowbreak{}real}}
\leansource{\texttt{wasserstein2\_\allowbreak{}le\_\allowbreak{}random\_\allowbreak{}cost\_\allowbreak{}real} \quad UnadjustedLangevin/\allowbreak{}Coupling.lean, lines 103–108}
\begin{leancode}
\leanline{0}{theorem\ wasserstein2\_le\_random\_cost\_real\ (X\ Y\ :\ Ω\ →\ H)}
\leanline{4}{(hX\ :\ Measurable\ X)\ (hY\ :\ Measurable\ Y)}
\leanline{4}{(hXY\ :\ MemLp\ (fun\ ω\ ↦\ X\ ω\ -\ Y\ ω)\ 2\ P)\ :}
\leanline{4}{wasserstein2Real\ (P.map\ X)\ (P.map\ Y)\ ≤}
\leanline{6}{(eLpNorm\ (fun\ ω\ ↦\ X\ ω\ -\ Y\ ω)\ 2\ P).toReal}
\end{leancode}
\leanlabel{How the Lean proof follows the argument}
In \texttt{wasserstein2\_\allowbreak{}le\_\allowbreak{}random\_\allowbreak{}cost\_\allowbreak{}real} the hypothesis \texttt{h\allowbreak{}XY : Mem\allowbreak{}Lp (fun \ensuremath{\omega} \ensuremath{\mapsto} X \ensuremath{\omega} -\allowbreak{} Y \ensuremath{\omega}) 2 P} is square integrability of the difference; \texttt{h\allowbreak{}XY.\allowbreak{}e\allowbreak{}Lp\allowbreak{}Norm\_\allowbreak{}ne\_\allowbreak{}top} is finiteness of its norm. Mathlib's \texttt{ENNReal.\allowbreak{}to\allowbreak{}Real\_\allowbreak{}mono} (if \(b\ne\infty\) and \(a\le b\) then \(a\) and \(b\) compare in the same way as real numbers) converts the inequality of Step 9. The right-hand side of the conclusion is \texttt{(e\allowbreak{}Lp\allowbreak{}Norm (fun \ensuremath{\omega} \ensuremath{\mapsto} X \ensuremath{\omega} -\allowbreak{} Y \ensuremath{\omega}) 2 P).\allowbreak{}to\allowbreak{}Real}, the finite value of \(\|X-Y\|_{2,P}\).
\end{leanbox}

\section{\texorpdfstring{The transport bound for classes in \(L^2\)}{The transport bound for classes in L²}}


\textbf{Claim.} Let \(P\) be a measure and \(X,Y\in L^2(P;H)\), with representatives also denoted \(X,Y\). Then


\[W_2(X_\#P,Y_\#P)\le\|X-Y\|_{L^2}<\infty.\]


\textbf{Proof.} Representatives of classes in \(L^2\) are almost everywhere strongly measurable, hence almost everywhere measurable, so Step 10 gives \(W_2(X_\#P,Y_\#P)\le\|\omega\mapsto X(\omega)-Y(\omega)\|_{2,P}\). The function \(\omega\mapsto X(\omega)-Y(\omega)\) agrees almost everywhere with the representative of the class \(X-Y\); hence its extended norm is \(\|X-Y\|_{L^2}\), which is finite. \(\square\)

\begin{leanbox}{Lean: \texttt{wasserstein2\_\allowbreak{}le\_\allowbreak{}Lp\_\allowbreak{}cost}}
\leansource{\texttt{wasserstein2\_\allowbreak{}le\_\allowbreak{}Lp\_\allowbreak{}cost} \quad UnadjustedLangevin/\allowbreak{}Coupling.lean, lines 112–119}
\begin{leancode}
\leanline{0}{theorem\ wasserstein2\_le\_Lp\_cost\ (X\ Y\ :\ Lp\ H\ 2\ P)\ :}
\leanline{4}{wasserstein2\ (P.map\ X)\ (P.map\ Y)\ ≤\ ENNReal.ofReal\ ‖X\ -\ Y‖}
\end{leancode}
\leanlabel{How the Lean proof follows the argument}
The conclusion of \texttt{wasserstein2\_\allowbreak{}le\_\allowbreak{}Lp\_\allowbreak{}cost} is stated in \texttt{\ensuremath{\mathbb R}\ensuremath{\ge}0\ensuremath{\infty}}, with the real norm embedded by \texttt{ENNReal.\allowbreak{}of\allowbreak{}Real}. \texttt{hcost} identifies the two extended norms through \texttt{e\allowbreak{}Lp\allowbreak{}Norm\_\allowbreak{}congr\_\allowbreak{}ae} and the almost-everywhere identity \texttt{Lp.\allowbreak{}coe\allowbreak{}Fn\_\allowbreak{}sub X Y}; \texttt{Lp.\allowbreak{}norm\_\allowbreak{}def} writes \(\|X-Y\|\) as the real value of the extended norm, and \texttt{ENNReal.\allowbreak{}of\allowbreak{}Real\_\allowbreak{}to\allowbreak{}Real} (with finiteness \texttt{Lp.\allowbreak{}e\allowbreak{}Lp\allowbreak{}Norm\_\allowbreak{}ne\_\allowbreak{}top}) undoes the conversion. The bound itself is \texttt{wasserstein2\_\allowbreak{}le\_\allowbreak{}ae\_\allowbreak{}random\_\allowbreak{}cost}, with almost everywhere measurability from \texttt{Lp.\allowbreak{}aestrongly\allowbreak{}Measurable}.
\end{leanbox}

\section{\texorpdfstring{The real transport bound for classes in \(L^2\)}{The real transport bound for classes in L²}}


\textbf{Claim.} For \(X,Y\in L^2(P;H)\),


\[W_2^{\mathbb R}(X_\#P,Y_\#P)\le\|X-Y\|_{L^2}.\]


\textbf{Proof.} By the previous step \(W_2(X_\#P,Y_\#P)\) is finite and at most \(\|X-Y\|_{L^2}\); hence \(W_2^{\mathbb R}(X_\#P,Y_\#P)=W_2(X_\#P,Y_\#P)\le\|X-Y\|_{L^2}\). \(\square\)




This is the form used at the end of the convergence proofs: if \(X\) is the output of the algorithm and \(Y\) has the target law \(\pi\), any bound on \(\|X-Y\|_{L^2}\) bounds \(W_2^{\mathbb R}(X_\#P,\pi)\).

\begin{leanbox}{Lean: \texttt{wasserstein2\_\allowbreak{}le\_\allowbreak{}Lp\_\allowbreak{}cost\_\allowbreak{}real}}
\leansource{\texttt{wasserstein2\_\allowbreak{}le\_\allowbreak{}Lp\_\allowbreak{}cost\_\allowbreak{}real} \quad UnadjustedLangevin/\allowbreak{}Coupling.lean, lines 122–126}
\begin{leancode}
\leanline{0}{theorem\ wasserstein2\_le\_Lp\_cost\_real\ (X\ Y\ :\ Lp\ H\ 2\ P)\ :}
\leanline{4}{wasserstein2Real\ (P.map\ X)\ (P.map\ Y)\ ≤\ ‖X\ -\ Y‖}
\end{leancode}
\leanlabel{How the Lean proof follows the argument}
\texttt{wasserstein2\_\allowbreak{}le\_\allowbreak{}Lp\_\allowbreak{}cost\_\allowbreak{}real} applies \texttt{ENNReal.\allowbreak{}to\allowbreak{}Real\_\allowbreak{}mono} to the previous step (the right-hand side \texttt{ENNReal.\allowbreak{}of\allowbreak{}Real \ensuremath{\Vert}X -\allowbreak{} Y\ensuremath{\Vert}} is finite by \texttt{ENNReal.\allowbreak{}of\allowbreak{}Real\_\allowbreak{}ne\_\allowbreak{}top}) and simplifies with \texttt{ENNReal.\allowbreak{}to\allowbreak{}Real\_\allowbreak{}of\allowbreak{}Real} and \texttt{norm\_\allowbreak{}nonneg}.
\end{leanbox}

\section{\texorpdfstring{Pointwise domination passes to the \(L^2\)-norm}{Pointwise domination passes to the L²-norm}}


\textbf{Claim.} Let \(P\) be a measure on \(\Omega\), \(A,D:\Omega\to H\) and \(q\ge0\). Assume that \(A\) is almost everywhere strongly measurable and that \(\|A(\omega)\|\le q\|D(\omega)\|\) for \(P\)-almost every \(\omega\). Then


\[\|A\|_{2,P}\le q\,\|D\|_{2,P}\qquad\text{in }[0,\infty].\]


\textbf{Proof.} If \(D\) is almost everywhere strongly measurable, squaring the pointwise bound (both sides are nonnegative) and integrating gives \(\int\|A\|^2\,dP\le q^2\int\|D\|^2\,dP\) in \([0,\infty]\); taking square roots gives the claim. Otherwise \(\|D\|_{2,P}=\infty\): if \(q>0\) the right-hand side is \(\infty\); if \(q=0\), then \(A=0\) almost everywhere, so \(\|A\|_{2,P}=0=0\cdot\infty\). \(\square\)




No relation between the laws of \(A\) and \(D\) is needed.

\begin{leanbox}{Lean: \texttt{eLpNorm\_\allowbreak{}contraction}}
\leansource{\texttt{eLpNorm\_\allowbreak{}contraction} \quad UnadjustedLangevin/\allowbreak{}Coupling.lean, lines 130–136}
\begin{leancode}
\leanline{0}{theorem\ eLpNorm\_contraction\ (A\ D\ :\ Ω\ →\ H)\ (q\ :\ ℝ≥0)}
\leanline{4}{(hA\ :\ AEStronglyMeasurable\ A\ P)}
\leanline{4}{(hcontract\ :\ ∀ᵐ\ ω\ ∂P,\ ‖A\ ω‖\ ≤\ (q\ :\ ℝ)\ *\ ‖D\ ω‖)\ :}
\leanline{4}{eLpNorm\ A\ 2\ P\ ≤\ (q\ :\ ℝ≥0∞)\ *\ eLpNorm\ D\ 2\ P}
\end{leancode}
\leanlabel{How the Lean proof follows the argument}
In \texttt{e\allowbreak{}Lp\allowbreak{}Norm\_\allowbreak{}contraction}, \texttt{q : \ensuremath{\mathbb R}\ensuremath{\ge}0}, \texttt{h\allowbreak{}A : AEStrongly\allowbreak{}Measurable A P}, and \texttt{hcontract : \ensuremath{\forall}ᵐ \ensuremath{\omega} \ensuremath{\partial}P, \ensuremath{\Vert}A \ensuremath{\omega}\ensuremath{\Vert} \ensuremath{\le} (q : \ensuremath{\mathbb R}) * \ensuremath{\Vert}D \ensuremath{\omega}\ensuremath{\Vert}}; the notation \texttt{\ensuremath{\forall}ᵐ \ensuremath{\omega} \ensuremath{\partial}P} means "for \(P\)-almost every \(\omega\)". The coercion \texttt{(q : \ensuremath{\mathbb R}\ensuremath{\ge}0\ensuremath{\infty})} embeds \(q\) into \([0,\infty]\). The proof is Mathlib's \texttt{e\allowbreak{}Lp\allowbreak{}Norm\_\allowbreak{}le\_\allowbreak{}nnreal\_\allowbreak{}smul\_\allowbreak{}e\allowbreak{}Lp\allowbreak{}Norm\_\allowbreak{}of\_\allowbreak{}ae\_\allowbreak{}le\_\allowbreak{}mul} at exponent \(2\); the pointwise hypothesis is converted to nonnegative-real norms by \texttt{exact\_\allowbreak{}mod\_\allowbreak{}cast}, and \texttt{ENNReal.\allowbreak{}smul\_\allowbreak{}def} rewrites the scalar multiplication as a product.
\end{leanbox}

\section{\texorpdfstring{One step: contraction plus defect}{One step: contraction plus defect}}


\textbf{Claim.} Let \(P\) be a measure on \(\Omega\), \(X,Y,X',Y',A,V:\Omega\to H\) and \(q\ge0\). Assume that \(A\) is almost everywhere strongly measurable and that, for \(P\)-almost every \(\omega\),


\[X'(\omega)-Y'(\omega)=A(\omega)-V(\omega),\qquad\|A(\omega)\|\le q\,\|X(\omega)-Y(\omega)\|.\]


Then


\[\|X'-Y'\|_{2,P}\le q\,\|X-Y\|_{2,P}+\|V\|_{2,P}\qquad\text{in }[0,\infty].\]


\textbf{Proof.} Since \(X'-Y'=A-V\) almost everywhere, \(\|X'-Y'\|_{2,P}=\|A-V\|_{2,P}\). By Minkowski's inequality, \(\|A-V\|_{2,P}\le\|A\|_{2,P}+\|V\|_{2,P}\), and by the previous step \(\|A\|_{2,P}\le q\|X-Y\|_{2,P}\). \(\square\)




In the application, \(X,Y\) are the current numerical and reference states, \(X',Y'\) the next ones, \(A\) the difference of the two updates with common noise, and \(V\) the local defect. The defect enters additively through Minkowski's inequality; no orthogonality between \(A\) and \(V\) is needed.

\begin{leanbox}{Lean: \texttt{synchronous\_\allowbreak{}l2\_\allowbreak{}bound}}
\leansource{\texttt{synchronous\_\allowbreak{}l2\_\allowbreak{}bound} \quad UnadjustedLangevin/\allowbreak{}Coupling.lean, lines 141–149}
\begin{leancode}
\leanline{0}{theorem\ synchronous\_l2\_bound\ (X\ Y\ X'\ Y'\ A\ V\ :\ Ω\ →\ H)\ (q\ :\ ℝ≥0)}
\leanline{4}{(hA\ :\ AEStronglyMeasurable\ A\ P)}
\leanline{4}{(hidentity\ :\ ∀ᵐ\ ω\ ∂P,\ X'\ ω\ -\ Y'\ ω\ =\ A\ ω\ -\ V\ ω)}
\leanline{4}{(hcontract\ :\ ∀ᵐ\ ω\ ∂P,\ ‖A\ ω‖\ ≤\ (q\ :\ ℝ)\ *\ ‖X\ ω\ -\ Y\ ω‖)\ :}
\leanline{4}{eLpNorm\ (fun\ ω\ ↦\ X'\ ω\ -\ Y'\ ω)\ 2\ P\ ≤}
\leanline{6}{(q\ :\ ℝ≥0∞)\ *\ eLpNorm\ (fun\ ω\ ↦\ X\ ω\ -\ Y\ ω)\ 2\ P\ +\ eLpNorm\ V\ 2\ P}
\end{leancode}
\leanlabel{How the Lean proof follows the argument}
\texttt{synchronous\_\allowbreak{}l2\_\allowbreak{}bound} rewrites the left-hand side with \texttt{e\allowbreak{}Lp\allowbreak{}Norm\_\allowbreak{}congr\_\allowbreak{}ae hidentity} (equal almost everywhere, equal norms), applies Minkowski's inequality \texttt{e\allowbreak{}Lp\allowbreak{}Norm\_\allowbreak{}sub\_\allowbreak{}le} for the exponent \(2\ge1\), and bounds the first term by \texttt{e\allowbreak{}Lp\allowbreak{}Norm\_\allowbreak{}contraction} (with \texttt{add\_\allowbreak{}le\_\allowbreak{}add} and \texttt{le\_\allowbreak{}rfl} for the second term).
\end{leanbox}

\section{\texorpdfstring{One step in transport distance}{One step in transport distance}}


\textbf{Claim.} Under the hypotheses of the previous step, assume in addition that \(X'\) and \(Y'\) are measurable. Then


\[W_2\bigl(X'_\#P,Y'_\#P\bigr)\le q\,\|X-Y\|_{2,P}+\|V\|_{2,P}.\]


\textbf{Proof.} By Step 9, \(W_2(X'_\#P,Y'_\#P)\le\|X'-Y'\|_{2,P}\); the previous step bounds the right-hand side. \(\square\)

\begin{leanbox}{Lean: \texttt{synchronous\_\allowbreak{}wasserstein\_\allowbreak{}bound}}
\leansource{\texttt{synchronous\_\allowbreak{}wasserstein\_\allowbreak{}bound} \quad UnadjustedLangevin/\allowbreak{}Coupling.lean, lines 152–160}
\begin{leancode}
\leanline{0}{theorem\ synchronous\_wasserstein\_bound\ (X\ Y\ X'\ Y'\ A\ V\ :\ Ω\ →\ H)\ (q\ :\ ℝ≥0)}
\leanline{4}{(hX'\ :\ Measurable\ X')\ (hY'\ :\ Measurable\ Y')}
\leanline{4}{(hA\ :\ AEStronglyMeasurable\ A\ P)}
\leanline{4}{(hidentity\ :\ ∀ᵐ\ ω\ ∂P,\ X'\ ω\ -\ Y'\ ω\ =\ A\ ω\ -\ V\ ω)}
\leanline{4}{(hcontract\ :\ ∀ᵐ\ ω\ ∂P,\ ‖A\ ω‖\ ≤\ (q\ :\ ℝ)\ *\ ‖X\ ω\ -\ Y\ ω‖)\ :}
\leanline{4}{wasserstein2\ (P.map\ X')\ (P.map\ Y')\ ≤}
\leanline{6}{(q\ :\ ℝ≥0∞)\ *\ eLpNorm\ (fun\ ω\ ↦\ X\ ω\ -\ Y\ ω)\ 2\ P\ +\ eLpNorm\ V\ 2\ P}
\end{leancode}
\leanlabel{How the Lean proof follows the argument}
\texttt{synchronous\_\allowbreak{}wasserstein\_\allowbreak{}bound} is \texttt{(wasserstein2\_\allowbreak{}le\_\allowbreak{}random\_\allowbreak{}cost P X' Y' h\allowbreak{}X' h\allowbreak{}Y').\allowbreak{}trans} applied to \texttt{synchronous\_\allowbreak{}l2\_\allowbreak{}bound}.
\end{leanbox}

\section{\texorpdfstring{The one-step inequality in real numbers}{The one-step inequality in real numbers}}


\textbf{Claim.} Under the hypotheses of Step 15, assume in addition that \(X-Y\) and \(V\) are square integrable. Then all quantities below are finite and


\[\|X'-Y'\|_{2,P}\le q\,\|X-Y\|_{2,P}+\|V\|_{2,P}\qquad\text{in }\mathbb R.\]


\textbf{Proof.} Since \(q<\infty\), \(\|X-Y\|_{2,P}<\infty\) and \(\|V\|_{2,P}<\infty\), the right-hand side of Step 15 is finite; hence so is the left-hand side, and the inequality of Step 15 is an inequality between real numbers. For finite values, the real value of a sum or product in \([0,\infty]\) is the sum or product of the real values. \(\square\)




In particular \(X'-Y'\) is square integrable whenever \(X-Y\) and \(V\) are. Iterating this one-step inequality gives the scalar recurrence \(e_{k+1}\le qe_k+b\) for \(e_k=\|X_k-Y_k\|_{2,P}\).

\begin{leanbox}{Lean: \texttt{synchronous\_\allowbreak{}l2\_\allowbreak{}bound\_\allowbreak{}real}}
\leansource{\texttt{synchronous\_\allowbreak{}l2\_\allowbreak{}bound\_\allowbreak{}real} \quad UnadjustedLangevin/\allowbreak{}Coupling.lean, lines 163–176}
\begin{leancode}
\leanline{0}{theorem\ synchronous\_l2\_bound\_real\ (X\ Y\ X'\ Y'\ A\ V\ :\ Ω\ →\ H)\ (q\ :\ ℝ≥0)}
\leanline{4}{(hA\ :\ AEStronglyMeasurable\ A\ P)}
\leanline{4}{(hidentity\ :\ ∀ᵐ\ ω\ ∂P,\ X'\ ω\ -\ Y'\ ω\ =\ A\ ω\ -\ V\ ω)}
\leanline{4}{(hcontract\ :\ ∀ᵐ\ ω\ ∂P,\ ‖A\ ω‖\ ≤\ (q\ :\ ℝ)\ *\ ‖X\ ω\ -\ Y\ ω‖)}
\leanline{4}{(hXY\ :\ MemLp\ (fun\ ω\ ↦\ X\ ω\ -\ Y\ ω)\ 2\ P)\ (hV\ :\ MemLp\ V\ 2\ P)\ :}
\leanline{4}{(eLpNorm\ (fun\ ω\ ↦\ X'\ ω\ -\ Y'\ ω)\ 2\ P).toReal\ ≤}
\leanline{6}{(q\ :\ ℝ)\ *\ (eLpNorm\ (fun\ ω\ ↦\ X\ ω\ -\ Y\ ω)\ 2\ P).toReal\ +\ (eLpNorm\ V\ 2\ P).toReal}
\end{leancode}
\leanlabel{How the Lean proof follows the argument}
In \texttt{synchronous\_\allowbreak{}l2\_\allowbreak{}bound\_\allowbreak{}real}, \texttt{hprod} (via \texttt{ENNReal.\allowbreak{}mul\_\allowbreak{}ne\_\allowbreak{}top}) and the sum condition (via \texttt{ENNReal.\allowbreak{}add\_\allowbreak{}ne\_\allowbreak{}top}) prove finiteness of the right-hand side; \texttt{ENNReal.\allowbreak{}to\allowbreak{}Real\_\allowbreak{}mono} converts \texttt{synchronous\_\allowbreak{}l2\_\allowbreak{}bound} to real numbers, and \texttt{ENNReal.\allowbreak{}to\allowbreak{}Real\_\allowbreak{}add}, \texttt{ENNReal.\allowbreak{}to\allowbreak{}Real\_\allowbreak{}mul} and \texttt{ENNReal.\allowbreak{}coe\_\allowbreak{}to\allowbreak{}Real} distribute the conversion over the sum and the product. Each norm in the statement appears as \texttt{(e\allowbreak{}Lp\allowbreak{}Norm f 2 P).\allowbreak{}to\allowbreak{}Real}, the finite value of \(\|f\|_{2,P}\).
\end{leanbox}

\end{document}
